% C4 Architecture Visualization Documentation Suite | Original educational content
% Standalone LaTeX root: no external style file, no network, no shell escape.
\documentclass[10pt,letterpaper]{article}
\usepackage[T1]{fontenc}
\usepackage[utf8]{inputenc}
\usepackage{lmodern,microtype}
\usepackage[letterpaper,margin=0.82in,headheight=15pt,headsep=18pt,footskip=26pt]{geometry}
\usepackage[table]{xcolor}
\usepackage{booktabs,tabularx,array,longtable}
\usepackage{enumitem}
\usepackage{fancyhdr,lastpage}
\usepackage{titlesec}
\usepackage{hyperref,xurl}
\usepackage{needspace}
\usepackage{listings}
\usepackage{amssymb}
\definecolor{C4Navy}{HTML}{17364B}
\definecolor{C4Teal}{HTML}{137C85}
\definecolor{C4Gray}{HTML}{526775}
\definecolor{C4Light}{HTML}{EEF4F5}
\definecolor{C4Line}{HTML}{CAD8DE}
\color{C4Navy}
\hypersetup{colorlinks=true,linkcolor=C4Teal,urlcolor=C4Teal,citecolor=C4Teal}
\setlength{\parindent}{0pt}
\setlength{\parskip}{5pt}
\setlength{\emergencystretch}{2em}
\raggedbottom
\setlist[itemize]{leftmargin=1.4em,topsep=3pt,itemsep=3pt}
\setlist[enumerate]{leftmargin=1.6em,topsep=3pt,itemsep=3pt}
\titleformat{\section}{\large\bfseries\sffamily\color{C4Navy}}{\thesection}{0.6em}{}
\titleformat{\subsection}{\normalsize\bfseries\sffamily\color{C4Teal}}{\thesubsection}{0.6em}{}
\titlespacing*{\section}{0pt}{16pt}{6pt}
\titlespacing*{\subsection}{0pt}{9pt}{4pt}
\pagestyle{fancy}
\fancyhf{}
\fancyhead[L]{\scriptsize\sffamily C4 / ARCHITECTURE VISUALIZATION}
\fancyhead[R]{\scriptsize\sffamily REFERENCE / EDUCATION}
\fancyfoot[L]{\scriptsize\color{C4Gray}ORIGINAL ENGINEERING SUITE / 2026}
\fancyfoot[R]{\scriptsize\thepage\ / \pageref*{LastPage}}
\renewcommand{\headrulewidth}{0.4pt}
\newcolumntype{Y}{>{\raggedright\arraybackslash}X}
\newcommand{\RuleLine}{\par\noindent\color{C4Line}\rule{\linewidth}{0.45pt}\color{C4Navy}\par}
\newcommand{\WritingSpace}[1]{\par\vspace{2pt}\noindent\fcolorbox{C4Line}{white}{\parbox[t][#1][t]{0.94\linewidth}{\hspace{1pt}}}\par\vspace{7pt}}
\newcommand{\Field}[1]{\par\textbf{#1}\par\WritingSpace{0.57in}}
\newcommand{\Prompter}[1]{\par{\small\color{C4Teal}\textbf{Authoring guidance:} #1}\par}
\newcommand{\SourceNote}[1]{\par\smallskip\noindent\colorbox{C4Light}{\parbox{0.94\linewidth}{\footnotesize\textbf{PROVENANCE.} #1}}\par\medskip}
\lstset{basicstyle=\ttfamily\footnotesize,breaklines=true,columns=fullflexible,showstringspaces=false,frame=single,rulecolor=\color{C4Line}}
\newcommand{\DocStart}[3]{%
\begin{document}\begin{titlepage}
\vspace*{0.55in}
{\sffamily\bfseries\small\color{C4Teal}#1}\par\vspace{0.24in}
{\sffamily\bfseries\fontsize{26}{31}\selectfont\color{C4Navy}#2}\par\vspace{0.18in}
{\large\color{C4Gray}#3}\par\vspace{0.4in}
\textcolor{C4Teal}{\rule{\linewidth}{1.15pt}}\par\vspace{0.24in}
\textbf{Document control}\par
Edition 1.0\quad|\quad 30 October 2026\quad|\quad Original educational engineering material\par
\vspace{0.17in}
\textbf{Source boundary}\par
Simon Brown, \emph{The C4 Model: Visualizing Software Architecture} (O'Reilly, 2026), is the primary study source. The suite's exercises, forms, crosswalks, process rules, sample systems, and scoring models are independently authored supplements; they are not represented as the author's prescriptions.\par
\vfill
{\small\color{C4Gray}Not an ISO/IEC/IEEE 42010 conformance certificate. Not a substitute for the published standard or for review of deployed systems.}\par
\end{titlepage}
}
\DocStart{TECHNICAL REFERENCE}{C4 Foundations and Abstraction Semantics}{book-derived interpretation}
\tableofcontents\clearpage
\SourceNote{Sections explicitly marked book-aligned reflect the uploaded C4 text at a conceptual level. Integrations, traceability and enforcement rules are original companion engineering guidance.}
\section{Architectural object model}
Person, software system, container, component and code form four nested technical zoom levels plus the interacting human roles. A container means an application or datastore, not a cloud resource container or Docker packaging. A component is an internal grouping inside one container. A code element is an implementation unit.


\subsection{Review questions}
\begin{itemize}
\item What stakeholder or design question does this section answer?
\item How would a reviewer confirm the claim against code, runtime, or other primary evidence?
\item Which exception needs explicit documentation rather than an invented assumption?
\end{itemize}

\section{Model boundaries and granularity}
Choose the entity of interest before naming elements. A context view contains only the central software system and adjacent persons/systems. A container view opens one boundary. A component view opens one container. A code view is a selective implementation zoom.


\subsection{Review questions}
\begin{itemize}
\item What stakeholder or design question does this section answer?
\item How would a reviewer confirm the claim against code, runtime, or other primary evidence?
\item Which exception needs explicit documentation rather than an invented assumption?
\end{itemize}

\section{Relationships as architecture evidence}
Every interaction should have a direction and verb phrase; label the protocol, technology or data where known and relevant. A static relationship means a possible dependency; a dynamic view shows one concrete scenario.


\subsection{Review questions}
\begin{itemize}
\item What stakeholder or design question does this section answer?
\item How would a reviewer confirm the claim against code, runtime, or other primary evidence?
\item Which exception needs explicit documentation rather than an invented assumption?
\end{itemize}

\section{Seven diagrams}
Four static zooms: context, container, component, code. Three supporting forms: dynamic scenario, deployment instances, and system landscape. Diagram selection is concern-led; the C4 method does not oblige all seven.


\subsection{Review questions}
\begin{itemize}
\item What stakeholder or design question does this section answer?
\item How would a reviewer confirm the claim against code, runtime, or other primary evidence?
\item Which exception needs explicit documentation rather than an invented assumption?
\end{itemize}

\section{Visual semantics}
Titles, visible type labels, self-explanatory responsibilities, key/legend, boundary labels, intentional arrow direction and meaningful technology details are more important than elaborate icons.


\subsection{Review questions}
\begin{itemize}
\item What stakeholder or design question does this section answer?
\item How would a reviewer confirm the claim against code, runtime, or other primary evidence?
\item Which exception needs explicit documentation rather than an invented assumption?
\end{itemize}

\section{Model-code gap}
Maintain stable element IDs and relate model constructs to source paths, module owners, deployed workloads and operational evidence. These reconciliation mechanisms are original suite controls, not book-mandated records.


\subsection{Review questions}
\begin{itemize}
\item What stakeholder or design question does this section answer?
\item How would a reviewer confirm the claim against code, runtime, or other primary evidence?
\item Which exception needs explicit documentation rather than an invented assumption?
\end{itemize}

\section{Reference review protocol}
\begin{enumerate}
\item Declare type, scope, audience and concern for each model.
\item Identify and version source evidence for architectural assertions.
\item Reconcile cross-view correspondences and report gaps.
\item Separate proposed policy from externally normative requirements.
\item Record reviewer, findings, changes and disposition.
\end{enumerate}

\section{Source register}
\begin{itemize}
\item PRIMARY: Simon Brown, The C4 Model: Visualizing Software Architecture, O’Reilly Media, 2026, Chapters 1–12 and Appendix (user-supplied study copy).
\item AUTHORITATIVE C4: https://c4model.com/ — model abstractions, seven diagram types, official review checklist.
\item STANDARD SCOPE: ISO/IEC/IEEE 42010:2022, official catalogue https://www.iso.org/standard/74393.html; conformance requires the full standard.
\item UML: Object Management Group, UML 2.5.1, https://www.omg.org/spec/UML/2.5.1.
\item METHOD: Clements et al., Documenting Software Architectures: Views and Beyond, 2nd ed., Addison-Wesley, 2010/2011.
\item TOOLING: https://plantuml-stdlib.github.io/C4-PlantUML/ and repository tooling/plantuml/README.md.
\end{itemize}

\end{document}
